\documentclass[11pt]{article}
\usepackage{amsmath,amssymb}
\usepackage[a4paper,margin=25mm]{geometry}
\usepackage{xurl}
\title{Arithmetic cancellation in the mixed low row}
\author{}
\date{}
\begin{document}
\maketitle

\section{The mixed low row}

The arbitrary row Cauchy inequality uses only the norm of the completed
row. An improvement must estimate its correlation with the additive
column before taking absolute values. The identities below isolate the
complete principal correlation. The conductor arguments then bound
specified ranges of the incomplete sum.

Write $I$ for the squarefree completed column and $J$ for its cubic divisor
variable. With $A=IJ^3$, the upstream definitions give
\[
 \operatorname{summand}(\Psi,W,T,I,J)
 =\gamma_2(I)\overline{\alpha(A)}\Psi(A)
       \frac{q_J^{1/2}}{T^{1/2}}W(q_A/T).
\]
Here $\gamma_2(I)=\texttt{gaussTwo}(I)$ is the normalized cubic Gauss coefficient. Its modulus
is one on supported squarefree columns. All zero masks remain in $\Psi$.
The marked row multiplies this expression by $1_{D\mid A}$.
The additive polynomial has the exact normalization
\[
 A_{\sigma,m}(v)=Y^{-3/2}\sum_s c_\sigma(s)
       W_1(q_s/Y)(q_s/Y)^{-1+iv}g_s(-m).
\]
Write $a_{\sigma,v}(s)=c_\sigma(s)W_1(q_s/Y)(q_s/Y)^{-1+iv}$.

For a finite field $F$, a multiplicative character $\chi$, and an additive
character $\psi$, put $G_\chi(h)=\sum_x\chi(x)\psi(hx)$.
The upstream finite Gauss convolution proves the exact identity
\[
 \sum_{m\in F}\rho(m)G_\chi(-m)
 =\begin{cases}
   \rho(-1)^{-1}\tau(\rho)(|F|-1),&\chi=\rho,\\
   0,&\chi\ne\rho.
  \end{cases}
\]
The multiplicative characters are extended by zero at zero, including
the principal character. In the mixed row, the good prime exponent of
$\rho$ is $v_{\mathfrak p}(I)+3v_{\mathfrak p}(J)$.
An absent prime factor is the constant function one instead. For a source
prime of exponent one this distinction has no effect since $G_\chi(0)=0$.
At an arbitrary local dual frequency $k$, the exact transform is
\[
 \sum_{m\in F}\rho(m)G_\chi(-m)\psi(km)
   =\chi(-1)^{-1}\tau(\chi)G_{\rho\chi^{-1}}(k).
\]
This identity includes the zeros of $\rho\chi^{-1}$.
For an absent row factor, the transform is instead $|F|\chi(k)$.
These formulas retain the arithmetic dependence needed for a mixed
nonzero frequency estimate.

Suppose first that $s$ and $I$ are squarefree. At every good prime the
complete principal condition is
\[
 v_{\mathfrak p}(I)+3v_{\mathfrak p}(J)
        \equiv v_{\mathfrak p}(s)\pmod6.
\]
Since the two squarefree exponents belong to $\{0,1\}$, this is equivalent
to equality of those exponents and evenness of $v_{\mathfrak p}(J)$.
Thus the good prime resonance is exactly $I=s$, $J=b^2$, and $A=sb^6$.
The original source and completed masks make $s$ and $A$ prime to $S$.
By \texttt{calibrationForSet\_span}, the calibration modulus is exactly
$\prod_{\mathfrak p\in S}\mathfrak p$. Its local residue character is
nonprincipal at every prime of $S$. The exact
\texttt{rowCoefficient\_apply} formula shows that all other dependence on
$m$ is $\chi_A(m)$. The target character, reciprocity phase, and completed
correction are evaluated at $A$ and do not depend on $m$.
Thus CRT separates the calibration mean from every good prime mean.
For nonempty $S$ the calibration mean is zero. The upstream theorem
\texttt{calibration\_mixed\_fourier\_zero} at $H=0$ therefore gives
\[
 \mathcal M_{\rm full\ complete\ mean}=0.
\]
This covers all original fixed phases and requires no loss depending on
height. The higher powers in the Mellin phase are in the smooth row
profile and are not part of the complete arithmetic mean.

Higher source powers cannot contribute to the complete principal mean.
At a good prime of source exponent $a\ge2$, its Gauss transform is
supported on $\mathfrak p^{a-1}\mid m$. If $\mathfrak p\mid A$, the
completed character retains a zero mask on $\mathfrak p\mid m$, so the
product is zero pointwise. If $\mathfrak p\nmid A$, the completed row is
constant at this prime and the complete mean of the Gauss transform is
zero. This last assertion follows by summing its defining additive
character over $m$. Every nonzero source residue has zero additive mean.

The good prime principal factor, before its zero calibration mean is
applied, has a bound stronger than the existing product of norm bounds.
Assume fixed compact windows $q_s\asymp Y$, $q_A\asymp T$, with
$T\ge cY>0$. On the resonance $J=b^2$ the completed weight is
$q_b/T^{1/2}$. The ideal count gives
\[
 \sum_{q_b\ll(T/Y)^{1/6}}\frac{q_b}{T^{1/2}}
 \ll T^{-1/6}Y^{-1/3}.
\]
The normalized local complete Gauss mean has modulus at most $q_s^{1/2}$.
There are $O(Y)$ source ideals. Consequently the zero frequency term of
an uncalibrated row Poisson mean, whose radial row integral is $O(Q)$,
satisfies
\[
 |\mathcal M_{\rm principal}|
 \ll QY^{-3/2}\,Y\,Y^{1/2}\,T^{-1/6}Y^{-1/3}
 =\frac Q{\sqrt Y}\left(\frac YT\right)^{1/6}.
\]
For the actual sum of surviving marked slots, every selected prime must
divide $A=sb^6$. Distinct slot supports imply at most
$\omega(A)^K\ll_{K,\epsilon}q_A^\epsilon$ tuples. Bounded slot weights
therefore add only $T^\epsilon$, with no factor from their total lengths.
The Gaussian annuli can be summed using their homogeneous seminorm decay.

There is a stronger bound for the actual resonant partial sum.
On $I=s$, $J=b^2$, the squarefree source Gauss identity gives
\[
 g_s(-m)\chi_{sb^6}(m)
   =\chi_s(-1)^{-1}\tau_s\,
       1_{(m,sb)=1}.
\]
The zero extension at primes dividing $b$ is retained.
After multiplication by the calibration character this coefficient has
zero mean and period norm
\[
 L_{s,b}\ll_{\rm fixed}q_s\operatorname{radnorm}(b)
       \ll_{\rm fixed}Y(T/Y)^{1/6}.
\]
For any fixed $a>2$, Poisson applied to the common compact smooth row
profile gives, whenever $Q\ge C Y(T/Y)^{1/6}$,
\[
 |\mathcal M_{\rm resonant}|
 \ll_{a,\epsilon}(1+|v|)^{J_a}
       Q T^{-1/6}Y^{-1/3}T^\epsilon
       \left(\frac{Y(T/Y)^{1/6}}Q\right)^a.
\]
Indeed its normalized arithmetic Fourier coefficients are bounded by
$q_s^{1/2}$. Its zero frequency vanishes by calibration. The other dual
vectors have linear scale $(Q/L_{s,b})^{1/2}$, and Fourier decay of
order $2a$ bounds their total by $(L_{s,b}/Q)^a$.
Summing the same positive source and divisor counts proves the display.
The exponent of the moving period is $(5l_y+N)/6$ when $Y=Z^{l_y}$
and $T=Z^N$. At the original empty subset geometry this is
$19/32$, below the row exponent $5/6$.
The estimate concerns the arithmetic resonant partial sum and uses the
original calibration. It does not change the principal Mellin signal of
the detector.

\subsection{A range of small residual conductors}

Fix $S$ and the calibration, all original ray classes, and the target
character. Let $f_v(z)=\Omega(|z|^2)|z|^{-2iv}$, with $\Omega$ supported
in a fixed positive compact annulus. Fix compact source windows at
$q_s\asymp Y$, $q_I\asymp C$, $q_J\asymp N$, where $CN^3\asymp T$.
The actual smooth row profile is $h_{\sigma,v}=f_v$ for every ray
$\sigma$. Its Fourier seminorms of any fixed order are bounded by a
fixed polynomial in $1+|v|$. All dependence on $\sigma$ is retained
in the arithmetic coefficients.
Assume $Y\ll_{\rm fixed}T$, as in the actual probe geometry.
The constants below are uniform in $Y,C,N,Q\ge1$, in the finite ray class
$\sigma$, in the Mellin height $v$, and in every real inverse row height $t$.
The completed window is
$W_t(x)=\texttt{logPhase}(t,\log x)W(x)$, and $|W_t(x)|=|W(x)|$.
It is evaluated at $q_A/T$ and has no row dependence. The absolute
estimates below do not differentiate it, so their constants are uniform
in $t$. They depend on the fixed windows
and their indicated seminorms. Further Fourier heights in the joint
source window can be included in the same finite height polynomial.

In this paragraph $s$ is squarefree. Write $R=\operatorname{rad}(sIJ)$.
Define two coprime squarefree ideals
\[
 H=\prod_{\substack{\mathfrak p\mid R\\
       v_{\mathfrak p}(I)+3v_{\mathfrak p}(J)
              -v_{\mathfrak p}(s)\not\equiv0\ (6)}}\mathfrak p,
 \qquad D=R/H.
\]
Let $\chi_H$ have the indicated nonzero sextic exponent at every prime
of $H$. It is primitive modulo $H$. The principal factors at primes of
$D$ remain the coprimality mask. The exact moving row coefficient is
\[
 g_s(-m)\chi_{IJ^3}(m)\xi(m)
    =\chi_s(-1)^{-1}\tau_s\chi_H(m)\xi(m)1_{(m,D)=1}.
\]
The finite fixed phases multiplying this identity have modulus bounded
independently of the moving ideals, by
\texttt{rowCoefficient\_apply} and the original completed correction
bound. They do not change $H$ or $D$.

For every $\epsilon>0$, the part with $q_H\le H_0$ satisfies the raw
mixed block estimate
\[
 |\mathcal M_{C,N;H\le H_0}|
   \ll_\epsilon (1+|v|)^J T^\epsilon
     \min\left\{\sqrt{CH_0},\frac{H_0}{\sqrt Y}\right\}.
\]
This assertion is absolute and does not use a lower bound for the norm
of the completed row.

To prove it, retain the exact mask and expand
$1_{(m,D)=1}=\sum_{d\mid D}\mu(d)1_{d\mid(m)}$.
The ideals $d$ are coprime to $H$ and to the calibration modulus.
Set $m=d n$, using its fixed primary generator. The row profile has norm
scale $Q/q_d$. The primitive character $\xi\chi_H$ has conductor norm
$\asymp_{\rm fixed}q_H$, and normalized Gauss coefficients of size
$O_{\rm fixed}(q_H^{-1/2})$. Smooth Poisson and Fourier decay give
\[
 \left|\sum_m\xi(m)\chi_H(m)1_{(m,D)=1}
                     f_v(m/\sqrt Q)\right|
 \ll_\epsilon (1+|v|)^J q_H^{1/2}T^\epsilon.
\]
For $Q_d=Q/q_d$, the Poisson majorant is a fixed multiple of
\[
 \frac{Q_d}{\sqrt{q_H}}\sum_{\ell\ne0}
       |\widehat f_v(\sqrt{Q_d/q_H}\,\ell)|.
\]
When $Q_d\le q_H$, lattice counting bounds this by $O(q_H^{1/2})$.
When $Q_d\ge q_H$, Fourier decay of order $2(a+1)$ bounds it by
$O(q_H^{1/2}(q_H/Q_d)^a)$.
Terms with $q_d$ larger than the fixed row radius are empty. The bounded
subunit scales cost a fixed constant. The divisor count supplies
$T^\epsilon$. When $Q/q_d\ge q_H$, the same argument with higher
Fourier decay supplies $(q_Hq_d/Q)^a$ for any fixed $a>2$.

Now put $g=(s,I)$. Every prime in $s/g$ or $I/g$ belongs to $H$, since
the squarefree exponents there differ and $3v_{\mathfrak p}(J)$ can be
only $0$ or $3$ modulo six. These two ideals are coprime. Hence
\[
 \frac{q_s q_I}{q_g^2}\le q_H\le H_0,
 \qquad q_g\gg\sqrt{YC/H_0}.
\]
On a gcd dyad $q_g\asymp G$, ideal counting bounds the number of pairs
$(s,I)$ by $O(YC/G)$. Summing the decreasing dyads from this lower
bound gives $O(\sqrt{YC H_0})$. The unrestricted pair count is also
$O(YC)$.
There are $O(N)$ divisor ideals $J$, with completed weight
$O(N^{1/2}/T^{1/2})$. The source Gauss factor has size $O(Y^{1/2})$.
Multiplication by $Y^{-3/2}$ therefore gives respectively
\[
 Y^{-1}\sqrt{H_0}\,\sqrt{YC H_0}\,
       \frac{N^{3/2}}{T^{1/2}}
       \asymp\frac{H_0}{\sqrt Y},
 \qquad
 Y^{-1}\sqrt{H_0}\,YC\,
       \frac{N^{3/2}}{T^{1/2}}
       \asymp\sqrt{CH_0}.
\]
The original summed marked slots add only $T^\epsilon$, because their
selected primes divide $A=IJ^3$. This proves the block bound.

For $H_0=Q/P^\delta$ with any fixed $\delta>0$, the small conductor
part is bounded by $(Q/\sqrt Y)P^{-\delta}T^\epsilon$.

There is also a bound for the whole block, including large conductors.
The radical bound gives $q_H\le q_R\ll YCN$.
Taking $H_0$ to be this fixed multiple of $YCN$ in the first branch gives
\[
 |\mathcal M_{C,N}|\ll_\epsilon(1+|v|)^J T^\epsilon C\sqrt{YN}
       \asymp(1+|v|)^J T^\epsilon\frac{T\sqrt Y}{N^{5/2}}.
\]
Consequently a factor $P^{-\delta}$ relative to $Q/\sqrt Y$ holds when
\[
 N\gg_{\rm fixed}(TY P^\delta/Q)^{2/5},
 \qquad\text{equivalently}\qquad
 C\ll_{\rm fixed}\left(\frac{Q}{YT^{1/6}P^\delta}\right)^{6/5}.
\]
At the original geometry, the thresholds for $\delta=0$ are
$N=Z^{13/40}$ and $C=Z^{23/120}$.
This estimate uses every actual marked weight through its divisor bound.
It requires no cancellation among completed columns and no norm lower bound.

The remaining squarefree-source sum has $q_H>Q/P^\delta$ and
$N\ll_{\rm fixed}(TY P^\delta/Q)^{2/5}$.
After the exact mask expansion
it is the sum of primitive character transforms against the coupled
$s,I,J$ weights
\begin{multline*}
 Y^{-3/2}T^{-1/2}\sum_{\substack{s\asymp Y,\ I\asymp C,\ J\asymp N\\
                   s,I\ \mathrm{squarefree},\ q_H>Q/P^\delta\\
                   N\ll(TYP^\delta/Q)^{2/5}}} \\
 a_{\sigma,v}(s)\gamma_2(I)\,q_J^{1/2}\chi_s(-1)^{-1}\tau_s
     \mathfrak a_\sigma(s,I,J)W(q_Iq_J^3/T)
     \mathcal R_{s,I,J},
\end{multline*}
where
\[
 \mathcal R_{s,I,J}=
     \sum_{d\mid D}\mu(d)\xi(d)\chi_H(d)
        \sum_n\xi(n)\chi_H(n)f_v(dn/\sqrt Q).
\]
Here $\mathfrak a_\sigma(s,I,J)$ retains the actual target, reciprocity,
calibration denominator, completed correction, marked prime sum, and
fixed phases. The source and completed windows are displayed explicitly.
No independent coefficient sequence has been substituted
for it. The two displayed block bounds do not control this remaining
range by $(Q/\sqrt Y)P^c$ with $c<1/12$.
It is the remaining sum in two coupled arithmetic variables for this route.

\subsection{Reduction of nonsquarefree sources}

Write every source uniquely as $s=t u$, where $t$ is squarefree, $u$ is
squareful, and $(t,u)=1$. Put
\[
 d=u/\operatorname{rad}(u),\qquad
 r=\prod_{\substack{\mathfrak p\mid u\\6\nmid v_{\mathfrak p}(u)}}\mathfrak p,
 \qquad
 e=\prod_{\substack{\mathfrak p\mid u\\6\mid v_{\mathfrak p}(u)}}\mathfrak p.
\]
Thus $u=d r e$. If $(u,A)\ne1$, the original mixed summand is zero
pointwise by the higher source power mask theorem.
Otherwise the exact prime power support enforces $d\mid(m)$.
After $m=d n$, CRT and the prime power lift identity give, up to a unit
phase independent of $n$,
\[
 g_s(-d n)
 =q_d\tau_{t,r}\chi_{t,r}(n)^{-1}
       \prod_{\mathfrak p\mid e}
          \bigl(q_{\mathfrak p}1_{\mathfrak p\mid(n)}-1\bigr),
 \qquad |\tau_{t,r}|=(q_tq_r)^{1/2}.
\]
The character $\chi_{t,r}$ has exponent one at primes of $t$ and
exponent $v_{\mathfrak p}(u)$ modulo six at primes of $r$.
It is primitive at each of these primes. The displayed principal factors
are exact Ramanujan sums, including their values at zero.

Expand the last product over $f\mid e$ with coefficient
$\mu(e/f)q_f$, and then substitute $n=f n'$. Since $(u,A)=1$, every
such $f$ is coprime to the active conductor and the calibration.
The active conductor is obtained from $t,I,J$ as before, with the
additional nonprincipal factors on $r$. Its norm satisfies
\[
 q_H\le q_tq_r\,C N\ll\frac{Y C N}{q_dq_e}.
\]
The passive masks remain exact. Primitive Poisson bounds the row after
both substitutions by $q_H^{1/2}(YT)^\epsilon$ uniformly in
$Q/(q_dq_f)$. Summing $f$ gives
\[
 q_d(q_tq_r)^{1/2}\sum_{f\mid e}q_f
     \ll_\epsilon (YT)^\epsilon\sqrt{Yq_dq_e}.
\]
More explicitly, after multiplying by the actual row character and
calibration, the row is
\[
 \begin{aligned}
 \eta_{s,A}q_d\tau_{t,r}
   &\sum_{f\mid e}\mu(e/f)q_f\xi(f)\chi_H(f) \\
   &\quad {}\times
   \sum_{h\mid D}\mu(h)\xi(h)\chi_H(h)
   \sum_k\xi(k)\chi_H(k)f_v(dfhk/\sqrt Q).
 \end{aligned}
\]
Here $\eta_{s,A}$ has modulus one and contains the source CRT phases,
the row factor $\chi_A(d)$, and $\xi(d)$.
The ideals $f,h$ are coprime to $H$ and the calibration. Their displayed
character values are therefore units. The last row has physical norm
scale $Q/(q_dq_fq_h)$. This form retains the original fixed phases and
both exact divisor expansions through every change of variables.
The product of this amplitude bound and $q_H^{1/2}$ is therefore at
most $Y\sqrt{C N}(YT)^\epsilon$, independently of $u$.
For fixed $u$, there are $O(Y/q_u)$ possible $t$. The same completed
column and cubic divisor counts now give
\[
 |\mathcal M_{C,N;\,u}|\ll_\epsilon
       (1+|v|)^J T^\epsilon\frac{C\sqrt{Y N}}{q_u}.
\]
The sum $\sum_{u\ {
m squareful}}q_u^{-1}$ converges. Its Euler factor
is $1+\sum_{a\ge2}q_{\mathfrak p}^{-a}=1+1/(q_{\mathfrak p}(q_{\mathfrak p}-1))$.
Thus the whole block estimate $C\sqrt{Y N}T^\epsilon$ and the proved
large $J$ range hold for all additive sources.

For $q_H\le H_0$, unrestricted counting instead gives
\[
 |\mathcal M_{C,N;\,u,H\le H_0}|\ll_\epsilon
       (1+|v|)^J T^\epsilon
       \frac{\sqrt{C H_0}}{\sqrt{q_uq_r}}.
\]
The series $\sum_{u\ {\rm squareful}}(q_uq_r)^{-1/2}$ converges.
At a prime with source exponent $a\ge2$, its local weight is
$q_{\mathfrak p}^{-(a+1_{6\nmid a})/2}$. The least exponent is $3/2$.
The $\sqrt{C H_0}$ branch therefore also extends to all additive sources.

The stronger $H_0/\sqrt Y$ branch does not follow by this argument.
Put $g=(t,I)$. Its conductor constraint only gives
$q_tq_Iq_r/q_g^2\le H_0$. The refined gcd count consequently yields
$H_0/(\sqrt Y q_r)$ for each fixed $u$.
The sum of $q_r^{-1}$ is not bounded by $Y^\epsilon$. For sources
$u=b^6$, one has $r=1$, and there are $\asymp Y^{1/6}$ such ideals
below $Y$. The Ramanujan substitutions and their reduced row scales
must be retained to treat this branch more sharply.
This is a precise obstruction to extending the refined small conductor
estimate by absolute values. The remaining nonsquarefree sources lie
in the same small $J$ range, but their exact Ramanujan factors must be
kept in the residual sum.

The principal and conductor statements above are absolute estimates.
The actual global zero frequency vanishes. Neither statement implies a bound relative
to $\|B^\varnothing\|_2$, which may be smaller than the available upper
bound $Q^{1/2}Z^\epsilon$. It does show that this principal resonance
cannot alone saturate the existing product of the proved norm bounds.
To improve the low exponent, it remains necessary to bound the nonzero
row Poisson frequencies of the actual mixed sum. Its moving modulus can
be much larger than the physical row length, so zero complete mean
alone gives no power saving. The final probe also has the outer factor
$Q^{-1/2}$ and the Mellin integral. The displayed bounds concern the raw
mixed row sum before this normalization.

The local algebra and the one plane zero mean Fourier tail reduction
are proved in \path{lean/GramMixedArithmetic.lean}.
The conductor block estimates above are analytic arguments and are not
claimed as Lean theorems.
The exact finite prime models and the squarefree valuation classification
can be reproduced with
\texttt{uv run research/gram\_mixed.py}.

\end{document}
